\subsection{The chains deployments actually ship}
\label{sec:prodchains}

The configuration costed so far is ours, and a reader is entitled to ask what
signers actually ship. We read the signatures out of \NcorpProdAssets{} files
signed by production software, \NcorpProdGens{} signing applications from
camera and phone makers, AI image services, editing tools and news
publishers, collected from Wikimedia Commons, the IPTC's media-provenance samples, public datasets and vendor posts, and out of the C2PA public test files. This is a convenience sample of what is publicly posted, not a census of signers. Every signature in a file is counted, including those of ingredients, because the file carries them. A chain is identified by its issuer, the leaf's issuer name and every certificate above the leaf, so a signer that issues a leaf per image counts once. Table~\ref{tab:corpus} gives the result by the application that wrote the file; a chain that reaches files of several groups, such as a Pixel camera chain in an image later edited with Microsoft's tools, appears in each row, and the \emph{Distinct} row counts every chain once; \NcorpMultiGroup{} of the \NcorpProdChains{} chains appear in more than one group. The supplement lists every chain and where the files came from (Section~\ref{sec:supp-corpus}); the released inventory gives every file's source URL, licence and SHA-256.

\begin{table*}[t]
\caption{Certificate chains in 317 production-signed files from 28 signing applications, and in the C2PA public test files, read against the C2PA trust list of 29~October 2026, a convenience sample of public files (Section~\ref{sec:supp-corpus} of the supplement lists the chains and sources, and the released inventory every file with its URL, licence and SHA-256). A chain is identified by its issuer: the leaf's issuer name and every certificate after the leaf, so leaves issued per image under one CA count once. Every signature in a file counts, the active manifest's and its ingredients', under the group of the application that wrote the file; a chain carried by files of several groups (Google's camera chains in files that Microsoft's tools wrote, for example) is counted in each, so the rows do not add up, and the \emph{Distinct} row counts each chain once. \emph{Sent}: certificates in \texttt{x5chain}. \emph{Root}: chains that carry a self-signed root. \emph{Listed}: chains whose path ends at a trust-list anchor. \emph{Issuer}: chains whose leaf was issued directly by an anchor, the condition for conformant leaf-only carriage. \emph{2.x}: files with a version-2 claim. \emph{33p}: signatures carrying \texttt{x5chain} under label~33 in the protected bucket. \emph{TST}: signatures with an RFC~3161 token, and its median size. \emph{OCSP}: signatures carrying revocation evidence.}
\label{tab:corpus}
\footnotesize
\setlength{\tabcolsep}{3.1pt}
\begin{tabular}{lrrrrlrrrrrrrr}
\toprule
 & & & \multicolumn{2}{c}{Chains} & & & & & & & \multicolumn{2}{c}{TST} & \\
\cmidrule(lr){4-5}\cmidrule(lr){12-13}
Signer & Files & Sigs & ES256 & PS256 & Sent & Root & Listed & Issuer & 2.x & 33p & n & B & OCSP\\
\midrule
Google & 105 & 141 & 5 & 0 & 2 & 0 & 5 & 2 & 105 & 141 & 141 & 2,015 & 43\\
OpenAI & 53 & 83 & 3 & 1 & 2, 3 & 1 & 3 & 0 & 51 & 83 & 31 & 5,259 & 23\\
Microsoft & 66 & 94 & 6 & 1 & 1, 2, 3 & 1 & 4 & 2 & 6 & 27 & 85 & 5,931 & 20\\
Adobe & 26 & 27 & 0 & 2 & 2 & 0 & 0 & 0 & 0 & 27 & 27 & 3,664 & 26\\
News publishers & 24 & 35 & 4 & 6 & 1, 2, 3 & 2 & 0 & 0 & 3 & 32 & 25 & 5,625 & 2\\
Other & 43 & 45 & 4 & 1 & 1, 2 & 0 & 1 & 0 & 3 & 45 & 43 & 5,828 & 0\\
\midrule
\emph{Distinct} & 317 & 425 & 15 & 8 & 1, 2, 3 & 4 & 9 & 2 & 168 & 355 & 352 & 3,663 & 114\\
\midrule
C2PA test files & 25 & 41 & 1 & 3 & 2, 3 & 1 & 0 & 0 & 0 & 0 & 39 & 5,940 & 0\\
\bottomrule
\end{tabular}
\end{table*}


\emph{The test files are not production output.} The C2PA public test files
hold \NcorpTestChains{}
distinct chains, almost all Adobe test certificates, every file has a
version-1 claim, and none reaches the current trust list. They show what a
conformance test looks like, not what signers emit.

\emph{ECDSA leads, and RSA is far from gone.} Of the
\NcorpProdChains{} distinct production chains, \NcorpProdEs{} sign with ES256
and \NcorpProdPs{} with PS256, and \NcorpProdSigEs{} of the
\NcorpProdSigs{} signatures are ES256. Google's phones and media services use
ES256 and OpenAI's image service uses both; RSA, mostly 2048-bit, persists in
Adobe's tools, Microsoft's image service and most newsrooms. Both are
therefore the baseline, and Table~\ref{tab:pq} reports growth against each.

\emph{Depth is three, and the root mostly stays out.} Signers carry the leaf and one intermediate in \NcorpProdTwoCert{} of the \NcorpProdChains{} chains; \NcorpProdRoot{} of the
\NcorpProdChains{} still carry the root, which the specification says to
leave out, and \NcorpProdOneCert{} carry the leaf alone, none of them with an
issuer on the trust list, so no validator using that list can take them to
\emph{Trusted}. The median chain carried is \NcorpProdChainMed{}
bytes.

\emph{The leaf-only condition already holds for some signers.}
Of the \NcorpProdChains{} chains, \NcorpProdListed{} end at an anchor on the
C2PA trust list, which was established in 2025 and does not yet cover many
earlier signers. Two of them, both Google's, are issued directly by a listed
anchor, so the path a validator checks is already one signature; they also
appear in the Microsoft row because Pixel images were edited there. Both still
carry that anchor beside the leaf, which the specification says not to
include, and dropping it would make them leaf-only and save about 730 bytes
each.

\emph{Timestamps and revocation are the ordinary case.} Of the
\NcorpProdSigs{} production signatures, \NcorpProdTs{} carry an RFC~3161 token, of
\NcorpProdTstLo{} to \NcorpProdTstHi{} bytes with a median of
\NcorpProdTstMed{}: about 2\,kB for an ECDSA time-stamping authority with two
certificates, as on Google's devices, and about 6\,kB for an RSA-4096 one with
three. Our ES256 and PS256 tokens, \NtokenBytesES{} and \NtokenBytesPS{}
bytes, sit at the small end of that range. \NcorpProdOcsp{} production
signatures each staple at least one OCSP response, so the OCSP column of
Table~\ref{tab:timestamp} describes what a quarter of today's signatures
already carry, not a hardened variant.

\emph{The header has moved to the RFC~9360 form.} Of the
\NcorpProdNxfive{} production signatures, \NcorpProdPthree{} carry \texttt{x5chain} under the
integer label~33 in the protected bucket, as C2PA~2.5 says generators shall;
the rest use the deprecated string label in the unprotected bucket, all but
\NcorpProdVtwoString{} of them in version-1 claims, which validators must
still accept.

\emph{What ML-DSA-65 would add to these files.} Because size is additive in
the components (Equation~\ref{eq:size}), the migration can be priced on the
files themselves. For each of the \NcpqFiles{} production files we kept
everything it carries, claims, assertions, ingredients, thumbnails and the
number of certificates each signer and each time-stamping authority sends, and
replaced every signature, carried certificate and timestamp token with our
measured ML-DSA-65 counterpart of the same role, so production names and
extensions give way to our profile. Three things make this an estimate rather
than a measurement of re-signed files, all of which widen it: the files are not
actually re-signed; a production certificate may differ from our profile by a
few hundred bytes, which enters the per-certificate total directly
(Section~\ref{sec:limitations}); and stapled OCSP responses, whose size the
files do not expose, are left classical, with a post-quantum upper bound given
separately (Table~\ref{tab:corpusstrata}). Percentages here are of the original
file. The median file is \NcpqFileMedMB\,MB; the median growth is
\NcpqAddMedkB\,kB, or \NcpqPctMed\% of a file, and the median
manifest store grows by a factor of \NcpqStoreGrowth. The median hides the
tail: files with several signatures, such as edited images that keep their
ingredients' credentials, grow by up to \NcpqAddMaxkB\,kB, and
\NcpqOverTen{} of the \NcpqFiles{} files would grow by more than a tenth
(90th percentile \NcpqPctPNinety\%). Of those, \NcpqOverTenTop{} are generated
images of about 200\,kB from one AI image service; among the other
\NcpqFilesOther{} files, \NcpqOverTenOther{} grow that much. Size decides the share: files under
500\,kB grow by a median \NcpqSmallMedPct\%, files over 5\,MB by
\NcpqBigMedPct\%, and the \NcpqThreeSigFiles{} files that keep three or more
signatures gain a median \NcpqThreeSigMedkB\,kB (Table~\ref{tab:corpusstrata}
and Figure~\ref{fig:corpusgrowth} in the supplement). Giving the
\NcpqOcspSigs{} stapled responses a full post-quantum delegated response
each, with nothing subtracted for the classical one, raises the median growth
of the \NcpqOcspFiles{} files that carry one from \NcpqOcspFilesMedkB{} to
\NcpqOcspFilesBoundMedkB\,kB, an upper bound. For the services that sign small
images at volume, the post-quantum credential is a double-digit share of every
file they serve. The released data list every file.
